\subsection{How Indicators Were Chosen}

An indicator earns its place on a dashboard by changing what somebody does. That criterion
is easy to state and routinely ignored, and the usual result is a wall of counts that is
read once out of curiosity and never again. Each figure reported here was therefore required
to answer a specific question that somebody actually has, and the indicators fall into three
groups according to whose question it is.

The first group answers \emph{is the extraction still correct}, and its audience is whoever
maintains the pipeline. These are regression instruments: they exist because a parser change
can break a rule silently, producing a node set that is smaller or emptier or differently
distributed while every individual record still looks plausible. A sampled manual check does
not catch that. A number that moves does.

The second group answers \emph{what is in this document}, and its audience is a reader
trying to understand the matrix as an object rather than to look one activity up. A person
consulting the DAM sees one activity at a time by construction, which makes the shape of the
whole invisible to them however carefully they read.

The third group answers \emph{what does this document imply about the institution}, and its
audience is governance and audit. These are the indicators of Section~\ref{sec:govkpi}, and
they are specified rather than delivered.

The rest of this section takes the delivered indicators in turn and states, for each, what it
measures, why it was chosen, and what it turned out to show. Table~\ref{tab:dashkpi}
summarises them first.

\begin{longtable}{|p{0.27\textwidth}|p{0.11\textwidth}|p{0.14\textwidth}|p{0.36\textwidth}|}
\caption{Indicators computed by the delivered analytics dashboard.}\label{tab:dashkpi}\\
\hline
\rowcolor{myblue}\textcolor{white}{\textbf{Indicator}} & \textcolor{white}{\textbf{Value}} & \textcolor{white}{\textbf{Answers}} & \textcolor{white}{\textbf{Why it is reported}} \\
\hline
\endfirsthead
\hline
\rowcolor{myblue}\textcolor{white}{\textbf{Indicator}} & \textcolor{white}{\textbf{Value}} & \textcolor{white}{\textbf{Answers}} & \textcolor{white}{\textbf{Why it is reported}} \\
\hline
\endhead
DAM nodes & 327 & What is in it & Establishes the size of the extracted corpus and moves immediately if a parsing change drops records. \\ \hline
Responsibilities extracted & 1{,}776 & What is in it & The real size of the matrix. Activities are containers; obligations are the content. \\ \hline
Distinct roles & 131 & What is in it & Bounds how many organisational actors the delegation structure involves at all. \\ \hline
Graph edges & 2{,}108 & What is in it & Confirms the graph is connected rather than a set of isolated records. \\ \hline
Unresolved role rate & 0.6\% & Is it correct & The single most sensitive regression detector in the project. \\ \hline
Average responsibilities per node & 5.48 & What is in it & Procedural weight per activity, and a second regression signal. \\ \hline
Nodes with no direct responsibilities & 29.6\% & What is in it & Quantifies the redirect structure, and justifies a correctness requirement of the agent. \\ \hline
Node type distribution & --- & What is in it & Separates structural scaffolding from actionable activity. \\ \hline
Authority code distribution & --- & What is in it & Shows which obligations the institution actually uses, against the seventeen available. \\ \hline
Top roles by responsibility count & --- & What is in it & Ranks actors by load and exposes concentration. \\ \hline
\end{longtable}

\subsection{Correctness Indicators}

\paragraph{Unresolved role rate.} This measures the proportion of responsibilities whose
authority code was located in a cell but whose owning column could not be determined. It
exists because of how role attribution works: column headers in the source are printed
rotated through ninety degrees, are not recoverable as table headers, and must be
reconstructed by clustering header coordinates and matching each code to a column by
horizontal alignment. That reconstruction is the most fragile step in the pipeline, and when
it degrades it does so quietly, attaching obligations to no role at all rather than raising
an error.

It was added as a development instrument rather than as an analytical one. During the
geometric attribution work it was the number watched after every change, because a
modification that improved one page's clustering while breaking another's is invisible in a
sampled check and unmistakable here. Its present value of 0.6\% is the figure against which
the pipeline's adequacy is argued in Chapter~\ref{ch:pipeline}, and the ten or so
responsibilities behind that percentage are the residue of genuinely ambiguous layout rather
than a systematic fault.

\paragraph{Average responsibilities per node.} At 5.48, this serves a similar purpose from
the opposite direction. The unresolved rate catches codes that lost their role; this catches
codes that were never found. A parsing change that stops recognising a code in some layout
leaves every remaining record perfectly valid, so nothing fails, and the only symptom is that
activities become quietly thinner. Watching a mean rather than a total matters here, because
a total also moves when the processed scope changes, and a ratio does not.

\subsection{Structural Indicators}

\paragraph{Node and responsibility counts.} The corpus contains 327 records carrying 1,776
individual responsibilities. Reporting both is deliberate, because the second is the real
measure of the matrix and the first is routinely mistaken for it. An activity is a container;
the obligations attached to it are the content, and a document with few activities and many
obligations each is a very different thing from the reverse.

\paragraph{Node type distribution.} Records are counted by type, which separates the document
into structural scaffolding and actionable content. Chapters and processes exist to organise;
tasks, child tasks and threshold variants are what a user actually asks about. The
distribution is what makes it possible to say honestly how much of a 79-page document is
answerable material rather than hierarchy, and it also exposes the threshold-variant
population, which is the mechanism behind the near-identical-activities problem that motivated
the whole project in Chapter~\ref{ch:context}.

\paragraph{Authority code distribution.} Seventeen level-qualified notations are available in
the legend. This indicator shows which of them the institution actually uses and how often,
and the interesting result is the imbalance: the delegation structure leans heavily on a
small subset, while several codes appear rarely enough that their meaning is unlikely to be
familiar to most staff. That is an argument for the glossary expansion built into the agent,
grounded in a measurement rather than an assumption.

\paragraph{Distinct roles and the role ranking.} There are 131 distinct roles, and the top
fifteen by obligation count are ranked. The count bounds the size of the delegation structure;
the ranking exposes its shape. A long flat ranking would indicate authority spread evenly, and
a steep one indicates concentration in a handful of offices. The delivered ranking is visibly
steep, which is the observation that motivates the authority-concentration indicator specified
in Section~\ref{sec:govkpi}.

\paragraph{Graph edges.} The 2,108 edges confirm that the extraction produced a connected
structure rather than a pile of independent records. It is reported because the difference
between a list and a graph is the entire premise of the approach: if containment and reference
edges were absent, the cross-reference following of Chapter~\ref{ch:agent} would have nothing
to traverse.

\subsection{An Indicator That Changed the Design}

\paragraph{Records with no direct responsibilities.} At 29.6\%, this is the figure with the
largest consequence in the chapter, and it was not anticipated.

It counts responsibility-bearing records that carry no obligation of their own. Encountering
it for the first time, the natural reading is that something is broken: nearly a third of the
extracted activities appear empty. Checking those records against the printed pages showed the
extraction was correct and the document is simply built that way. Roughly three activities in
ten record nothing themselves and point to another activity instead.

That reframed a feature as a requirement. Cross-reference following had been treated as a
refinement, worth having for completeness. This measurement showed that a system without it
dead-ends on close to a third of the matrix, which is not a gap in polish but a failure on
routine questions. The redirect path described in Chapter~\ref{ch:agent} exists in the form it
does because this number was measured before the agent was finished.

It is also the clearest demonstration of why the dashboard was worth building. The fact was
always present in the document and always invisible to a reader, because a reader
encountering one empty activity concludes that this activity is unusual. Only the aggregate
reveals that it is the norm.

